\documentclass[11pt]{article}
\usepackage{mathunal-base}
\mathunalfuente{serif}

\renewcommand{\examtitulo}{Cálculo en Varias Variables --- Tercer Parcial (Tema B)}
\renewcommand{\examfecha}{Semestre 01-2023, 21 de junio de 2023}

\begin{document}

\examheader

\vspace{4pt}
\noindent Nombre: \dotfill\ Cédula: \dotfill

\vspace{4pt}
\noindent Grupo: \dotfill\ Firma: \dotfill

\vspace{4pt}
\noindent Puntaje. Solo para uso oficial:\quad 1--4 (/52) \dotfill;\quad 5 (/24) \dotfill;\quad 6 (/24) \dotfill.\quad Total: \dotfill\quad Nota: \dotfill

\vspace{6pt}
\noindent\textbf{Instrucciones.} La duración del examen es de 1 hora y 50 minutos. El examen consta de seis puntos en dos hojas impresas por ambos lados; verifique que su examen esté completo. En las preguntas con procedimiento justifique sus respuestas en los espacios asignados. No está permitido sacar hojas en blanco ni ningún tipo de apuntes durante el examen; verifique que su celular esté apagado. No se permite el uso de calculadora.

\vspace{8pt}
\noindent\textbf{En las preguntas 1 a 4 no se tendrá en cuenta el procedimiento.}

\begin{enumerate}
\item{} (13\,\%) Sea $C$ la curva formada por la unión del semicírculo $D_1 = \left\{(x,y)\in\mathbb{R}^2 : x^2 + y^2 = 9,\ x \geq 0\right\}$ y el segmento que va desde $(0,-3)$ hasta $(0,3)$. Calcule
\[
\int_C \left(3y + \tan^{-1}(x^2)\right)dx + \left(2x^2 - \cos^{-1} y\right)dy.
\]

\vspace{2pt}
\noindent Respuesta: \underline{\phantom{XXXXXXXXXXXX}}

\item{} (13\,\%) Una partícula recorre la curva dada por $r(t) = (t,\ t^3,\ 2)$, con $t \in [0,2]$. Calcular el trabajo total realizado si el movimiento es producido por el campo de fuerzas
\[
\mathbf F(x,y,z) = e^x\,\vec i + xe^{3z}\,\vec j + x\cos(\pi y^2)\,\vec k.
\]

\vspace{2pt}
\noindent Respuesta: \underline{\phantom{XXXXXXXXXXXX}}

\item{} (13\,\%) Sea $\mathbf F$ el campo vectorial definido por $\mathbf F(x,y) = \left(\dfrac{1}{y},\ -\dfrac{x}{y^2}\right)$. Evalúe la integral
\[
\int_C \mathbf F \cdot dr
\]
si $C$ es la curva dada por $r(t) = (t\sin t,\ t^2\cos t)$, $\pi/4 \leq t \leq \pi/3$.

\vspace{2pt}
\noindent Respuesta: \underline{\phantom{XXXXXXXXXXXX}}

\item{} (13\,\%) Determine el área de la parte de la superficie $z = 4 - x^2 - y^2$ comprendida entre los planos $z = 1$ y $z = 3$.

\vspace{2pt}
\noindent Respuesta: \underline{\phantom{XXXXXXXXXXXX}}
\end{enumerate}

\vspace{4pt}
\noindent\textbf{En los ejercicios 5 y 6 debe realizar el procedimiento justificando cada paso.}

\begin{enumerate}[start=5]
\item{} (24\,\%) Sean $S$ la superficie dada por $S = \left\{(x,y,z) : y^2 + z^2 = 1,\ z \geq 0,\ 0 \leq x \leq 1\right\}$ y $C$ la frontera de $S$ orientada de modo antihorario vista desde el punto $(5,5,5)$. Encuentre el trabajo realizado por el campo de fuerzas
\[
\mathbf F(x,y,z) = (xy)\,\vec i + \left(xz^2(1-x)\right)\vec j + \left(x\sqrt{1-y^2}\right)\vec k,
\]
para mover una partícula a lo largo de $C$.

\vspace{2pt}
\noindent\footnotesize\textit{Nota de transcripción: la tercera componente del campo aparece parcialmente tachada en el examen original; se reconstruye aquí como $x\sqrt{1-y^2}$ a partir del desarrollo del estudiante.}\normalsize

\item{} (24\,\%) Considere la superficie formada por el cono $z = 4 - \sqrt{x^2+y^2}$ con $1 \leq z$. Use el teorema de la divergencia para evaluar
\[
\iint_S \mathbf F \cdot d\mathbf S,
\]
donde $\mathbf F(x,y,z) = (xy^2,\ x^2y,\ xy + 2)$ (suponga que el cono está orientado hacia arriba).
\end{enumerate}

\examfooter

\end{document}
